\documentclass[14pt, a4paper]{extarticle}
\usepackage{tz_style}

% Определение данных проекта
\newcommand{\developer}{ООО «Рога и копыта»}
\newcommand{\customer}{Московский Политех}
\newcommand{\systemtype}{Веб-приложение}
\newcommand{\systemname}{«Shortcuter»}
\newcommand{\systemfullname}{Веб-приложение для сокращения ссылок}
\newcommand{\teamlead}{Стальмахов Иван Сергеевич}
\newcommand{\group}{241-3211}
\newcommand{\teacher}{Кружалов Алексей Сергеевич}

\begin{document}

% Лист утверждения
\approvalpage{\developer}{\customer}{\systemtype}{\systemname}{14}{01.09.2026}

% Страница 2 — пустая (по ГОСТ 34.602-89 после титульного листа идёт пустая страница или аннотация)
\thispagestyle{empty}
\mbox{}
\newpage

% Содержание
\tableofcontents
\newpage

% 1. Общие положения
\section{Общие положения}

\subsection{Наименование системы}

Полное наименование системы -- \systemfullname.

Краткое наименование системы -- \systemname.

\subsection{Основания для проведения работ}

Работа выполняется в рамках учебного курса «Автоматизация процессов жизненного цикла программного обеспечения» в \customer.

\subsection{Наименование организаций -- Заказчика и Разработчика}

\textbf{Заказчик:} \customer. Адрес: 107023, г. Москва, ул. Большая Семёновская, д. 38. Телефон: 8 (495) 223-05-23.

\textbf{Разработчик:} \developer. Адрес: г. Москва. Телефон: +7 (495) 000-00-00.

\subsection{Перечень документов, на основании которых создается система}

Основанием для разработки \systemname является задание на лабораторную работу №1 по предмету «Автоматизация процессов жизненного цикла программного обеспечения».

\subsection{Плановые сроки начала и окончания работы}

Плановый срок начала работ: 4 сентября 2024 года.

Плановый срок окончания работ: 25 декабря 2024 года.

\subsection{Сведения об источниках и порядке финансирования работ}

Финансирование не предусмотрено (учебный проект).

\subsection{Порядок оформления и предъявления заказчику результатов работ}

Система передается в виде функционирующего веб-приложения, развернутого на тестовом хостинге. Приемка осуществляется преподавателем курса. Совместно с предъявлением системы сдается комплект документации согласно разделу 7 настоящего ТЗ.

\subsection{Состав используемой нормативно-технической документации}

При разработке системы и создании проектно-эксплуатационной документации используются следующие нормативные документы:

\begin{enumerate}
	\item ГОСТ 34.602-89. Техническое задание на создание автоматизированной системы.
	\item ГОСТ 34.601-90. Автоматизированные системы. Стадии создания.
	\item IEEE Std 830-1998. Рекомендации по составлению спецификации требований к программному обеспечению.
\end{enumerate}

% 2. Назначение и цели создания системы
\section{Назначение и цели создания системы}

\subsection{Назначение системы}

\systemname представляет собой веб-приложение для сокращения длинных URL-адресов. Система позволяет пользователям:

\begin{itemize}
	\item генерировать короткие ссылки из длинных URL;
	\item управлять созданными ссылками (просмотр, редактирование, удаление);
	\item отслеживать статистику переходов по ссылкам;
	\item настраивать пользовательские алиасы для ссылок.
\end{itemize}

Система предусматривает разграничение доступа по ролям пользователей: гость, зарегистрированный пользователь, администратор.

\subsection{Цели создания системы}

Основными целями разработки \systemname являются:

\begin{enumerate}
	\item создание удобного инструмента для сокращения URL-адресов;
	\item обеспечение управления ссылками с разграничением прав доступа;
	\item предоставление статистики использования ссылок;
	\item реализация системы с возможностью дальнейшего расширения функциональности.
\end{enumerate}

% 3. Характеристика объектов автоматизации
\section{Характеристика объектов автоматизации}

\subsection{Описание объекта автоматизации}

Объектом автоматизации является процесс создания, управления и отслеживания коротких ссылок. Система автоматизирует:

\begin{itemize}
	\item преобразование длинных URL в короткие идентификаторы;
	\item хранение и индексацию ссылок в базе данных;
	\item подсчет статистики переходов;
	\item управление пользовательскими ссылками через веб-интерфейс.
\end{itemize}

\subsection{Участники процесса}

В системе выделяются следующие группы пользователей:

\begin{table}[h]
\centering
\caption{Группы пользователей системы}
\label{tab:users}
\begin{tabular}{|p{3.5cm}|p{5cm}|p{5cm}|}
\hline
\textbf{Группа} & \textbf{Описание} & \textbf{Основные функции} \\
\hline
Гость & Неавторизованный пользователь & Создание короткой ссылки (без управления) \\
\hline
Пользователь & Зарегистрированный пользователь & Создание ссылок, управление своими ссылками, просмотр статистики \\
\hline
Администратор & Привилегированный пользователь & Полный доступ ко всем ссылкам, модерация, управление пользователями \\
\hline
\end{tabular}
\end{table}

% 4. Требования к системе
\section{Требования к системе}

\subsection{Требования к системе в целом}

\subsubsection{Требования к функциональности}

Система должна обеспечивать:

\begin{enumerate}
	\item создание короткой ссылки из произвольного URL;
	\item автоматическую генерацию уникального короткого идентификатора;
	\item возможность задания пользовательского алиаса (при доступности);
	\item редирект с короткой ссылки на оригинальный URL;
	\item подсчет количества переходов по каждой ссылке;
	\item хранение истории созданных ссылок для зарегистрированных пользователей;
	\item возможность редактирования и удаления своих ссылок;
	\item административный интерфейс для управления всеми ссылками и пользователями.
\end{enumerate}

\subsubsection{Требования к надежности}

Система должна:

\begin{itemize}
	\item обеспечивать сохранность данных при сбоях;
	\item предусматривать резервное копирование базы данных;
	\item корректно обрабатывать ошибочные ситуации (недоступность БД, невалидные URL).
\end{itemize}

\subsubsection{Требования к безопасности}

\begin{itemize}
	\item аутентификация пользователей по логину и паролю;
	\item разграничение доступа на уровне ролей;
	\item защита от несанкционированного доступа к чужим ссылкам;
	\item валидация всех входных данных.
\end{itemize}

\subsection{Требования к видам обеспечения}

\subsubsection{Требования к информационному обеспечению}

Хранение данных осуществляется в реляционной базе данных. Структура БД должна поддерживать:

\begin{itemize}
	\item таблицу пользователей;
	\item таблицу ссылок (связь с пользователем, оригинальный URL, короткий ID, дата создания);
	\item таблицу статистики переходов (связь со ссылкой, дата, IP, User-Agent).
\end{itemize}

\subsubsection{Требования к программному обеспечению}

\begin{itemize}
	\item Серверная часть: Python 3.11+ с фреймворком Flask/Django;
	\item База данных: PostgreSQL 14+;
	\item Клиентская часть: современный веб-браузер (Chrome, Firefox, Edge).
\end{itemize}

\subsubsection{Требования к техническому обеспечению}

Система должна работать на стандартном веб-хостинге с поддержкой Python и PostgreSQL. Минимальные требования:

\begin{itemize}
	\item 1 CPU, 512 MB RAM;
	\item 1 GB дискового пространства.
\end{itemize}

% 5. Состав и содержание работ по созданию системы
\section{Состав и содержание работ по созданию системы}

\begin{longtable}{|p{4cm}|p{3cm}|p{7cm}|}
\hline
\textbf{Стадия} & \textbf{Сроки} & \textbf{Этапы работ} \\
\hline
\endhead
1. Формирование требований & 01.09.2024 -- 15.09.2024 & Обследование объекта, формирование требований, оформление ТЗ \\
\hline
2. Разработка концепции & 01.09.2024 -- 15.09.2024 & Изучение аналогов, разработка архитектуры, выбор стека \\
\hline
3. Проектирование & 16.09.2024 -- 30.09.2024 & Проектирование БД, API, интерфейсов \\
\hline
4. Разработка & 01.10.2024 -- 30.11.2024 & Реализация backend, frontend, интеграция \\
\hline
5. Тестирование & 01.12.2024 -- 15.12.2024 & Unit-тесты, интеграционные тесты, исправление ошибок \\
\hline
6. Ввод в действие & 16.12.2024 -- 25.12.2024 & Деплой, подготовка документации, сдача \\
\hline
\end{longtable}

% 6. Порядок контроля и приемки системы
\section{Порядок контроля и приемки системы}

\subsection{Виды, состав, объем и методы испытаний}

Испытания системы проводятся в соответствии с программой и методикой испытаний, разрабатываемой в составе рабочей документации. Включают:

\begin{itemize}
	\item unit-тестирование отдельных модулей;
	\item интеграционное тестирование API;
	\item функциональное тестирование пользовательских сценариев.
\end{itemize}

\subsection{Общие требования к приемке работ}

Сдача-приемка работ производится поэтапно в соответствии с календарным планом. По результатам приемки подписывается акт сдачи-приемки работ.

% 7. Требования к документированию
\section{Требования к документированию}

Документирование проекта осуществляется в соответствии с требованиями ГОСТ 34.602-89. В состав документации входят:

\begin{enumerate}
	\item Техническое задание (настоящий документ);
	\item Руководство пользователя;
	\item Описание API (при наличии);
	\item Инструкция по развертыванию.
\end{enumerate}

Все документы должны быть предоставлены в электронном виде (PDF и исходные форматы).

% Список используемых источников
\section*{Список используемых источников}
\addcontentsline{toc}{section}{Список используемых источников}

\begin{enumerate}
	\item ГОСТ 34.602-89. Техническое задание на создание автоматизированной системы. -- М.: Стандартинформ, 1989.
	\item IEEE Std 830-1998. IEEE Recommended Practice for Software Requirements Specifications. -- IEEE, 1998.
	\item Официальный сайт Flask. URL: \url{https://flask.palletsprojects.com} [дата обращения: 11.09.2026].
	\item Официальный сайт PostgreSQL. URL: \url{https://www.postgresql.org} [дата обращения: 11.09.2026].
\end{enumerate}

% Подписи
\signatureblock{\developer}{Разработчик}

% Приложение с сокращениями
\abbreviations{
	URL & Uniform Resource Locator -- унифицированный указатель ресурса \\
	\hline
	API & Application Programming Interface -- интерфейс программирования приложений \\
	\hline
	БД & База данных \\
	\hline
	CRUD & Create, Read, Update, Delete -- базовые операции с данными \\
	\hline
	GUI & Graphical User Interface -- графический интерфейс пользователя \\
	\hline
	HTTP & HyperText Transfer Protocol -- протокол передачи гипертекста \\
	\hline
	ТЗ & Техническое задание \\
	\hline
}

\end{document}
